\begin{theorem}[Literal first rightward endpoint step]
\label{thm:peps_initial_string_right_step}
\lean{
TNLean.PEPS.regularWalkHolonomy_torusSweptStringRightStep,
TNLean.PEPS.torusSweptStringRightStepOperators,
TNLean.PEPS.torusSweptStringRightStepOperators_right_transport,
TNLean.PEPS.torusSweptStringRightStepOperators_up_transport
}
\leanok
On a torus with horizontal period at least eight and vertical period at least
seven, retain the two-string presentation with upward transport $g$ at four
consecutive columns and rightward transport $h^{-1}$ along its three-bond
partner string. Change only the middle upward transport at relative position
$(2,1)$ from $g$ to $gh$, with the inverse ordered coefficient at a periodic
seam. All rightward transports and all other upward transports stay literal.
The clockwise holonomies at relative bases $(1,1)$ and $(2,1)$ change from
$(h,1)$ to $(1,h)$.

This is the first elementary step of the prescribed four-endpoint route in
SCP10, Section~6.6 (local lines 2340--2423), within the stated finite-torus
scope. It does not establish the complete braid or reunion procedure;
see \cite{gap:rmp_peps_quantum_double_g_isometry}.
\uses{thm:peps_torus_swept_string_endpoint_holonomy,thm:peps_torus_plaquette_holonomy}
\end{theorem}
\begin{proof}\leanok
The native edge orientation gives the stated directed coefficients on either
seam. In the reverse plaquette product, the first changed holonomy is
$h(gh)^{-1}g=1$, whereas the second is $g^{-1}gh=h$.
\end{proof}

\begin{theorem}[Original-spin realization of the first rightward step]
\label{thm:peps_initial_string_right_physical_step}
\lean{
TNLean.PEPS.IsGIsometric.exists_unitary_torusInitialStringRightPhysicalStep,
TNLean.PEPS.torusFirstRightStepGauge
}
\leanok
Under the preceding period bounds, suppose the original site tensor is
$G$-isometric for the regular virtual action. One unitary on the six actual
spins in the horizontal two-plaquette block sends the initial internal
presentation to the preceding rightward-step presentation, for every $g,h$,
every virtual boundary configuration, and every common literal exterior and
crossing assignment. The unitary is chosen before all these parameters.
Its identity extension gives the corresponding equality of actual closed
contractions.

The continuing partner string is retained. This is a finite-torus elementary
physical movement from SCP10, Theorem~6.16 and the Section~6.6 route
(local lines 2271--2305 and 2340--2423), rather than a proof of the full braid.
\uses{thm:peps_initial_string_right_step,thm:peps_gauged_horizontal_flux_move}
\end{theorem}
\begin{proof}\leanok
On the lower row, assign vertex labels $h$ at the left vertex and the identity
at the other two vertices; assign $g$ throughout the upper row. The seven
actual internal bonds then reconstruct the initial presentation from a
single canonical $h$ insertion. Extending that insertion to the middle chord
changes precisely $g$ to $gh$ there. The common-background unitary acts on
all original-spin boundary columns. The shared exterior operators and the
actual cut expansion give the closed-contraction identity.
\end{proof}
